\documentclass[11pt]{article}

\usepackage[margin=1in]{geometry}
\usepackage[T1]{fontenc}
\usepackage[utf8]{inputenc}
\usepackage{lmodern}
\usepackage{microtype}
\usepackage{amsmath,amssymb,amsthm}
\usepackage{hyperref}
\usepackage{enumitem}

\title{Yang--Mills Grassroots Formalization Notes}
\author{Mettapedia autoformalization notes}
\date{\today}

\newcommand{\lean}[1]{\texttt{#1}}
\newcommand{\leansmall}[1]{{\small\texttt{#1}}}

\begin{document}
\maketitle

\section{Scope}

This note records first-party Lean work toward a Yang--Mills existence and mass
gap formalization.  The current target is deliberately modest: establish a
small spectral mass-gap audit layer that future classical-gauge and QFT
constructions must satisfy.

The Clay formulation says, in spectral terms, that the Hamiltonian has no
spectrum in an interval \((0,\Delta)\) above the vacuum.  The present Lean work
does not construct a Yang--Mills theory, a Hilbert space, a Hamiltonian, or the
Wightman/Osterwalder--Schrader infrastructure.  It only fixes the bounded
Hamiltonian gap predicates and proves basic consequences that should remain
valid under later strengthening.

\section{Lean Layer}

The first module is
\[
  \lean{Mettapedia/QuantumTheory/YangMills/MassGap.lean}.
\]
It introduces:
\begin{itemize}[leftmargin=1.5em]
  \item \lean{Spacetime} as \(\mathbb{R}^4\) and \lean{Space} as
    \(\mathbb{R}^3\), matching the coordinate types used by the existing
    reference statement;
  \item \lean{LinearOperator}, the bounded real-linear operator surface used by
    Mathlib's current spectrum API;
  \item \lean{SpectrumAvoidsOpenInterval}, a direct predicate for absence of
    spectrum in an interval;
  \item \lean{HasSpectralMassGap}, the positive spectral window
    \(0<\Delta\) plus absence of spectrum in \((0,\Delta)\);
  \item \lean{HasPositiveEnergySpectrum} and \lean{HasVacuumEnergyZero} as
    separate spectral audit assumptions;
  \item \lean{HasFirstPositiveSpectralValue}, the bounded-operator audit
    predicate for a first positive spectral excitation \(m\);
  \item \lean{PositiveSpectrumSet} and \lean{PositiveSpectralInf} as
    positive-spectrum order audit surfaces;
  \item \lean{SpectralGapSet}, \lean{SpectralGapSup}, and
    \lean{HasFiniteSpectralMass} as finite-gap audit surfaces;
  \item \lean{HasQuadraticMassGap}, the Hilbert-space quadratic-form version
    often used as a proof obligation before invoking spectral machinery.
\end{itemize}

\section{Proved Consequences}

The module proves that spectral gaps are downward closed:
\[
  \lean{HasSpectralMassGap.mono}.
\]
Thus any proposed proof that produces a gap \(\Delta\) also produces every
smaller positive gap.  This is a useful sanity check for future gap claims and
for any definition of the mass as a supremum of admissible gaps.

It also proves a pointwise obstruction:
\[
  \lean{not\_hasSpectralMassGap\_of\_spectrum\_in\_gap}.
\]
A single spectral value \(E\) with \(0<E<\Delta\) rules out \(\Delta\) as a
valid spectral mass gap.  This is the Yang--Mills analogue of a concrete audit
target: any future construction must explain why the Hamiltonian has no such
low positive spectral values.

The same exclusion is now available in threshold and set-subset forms:
\begin{center}
\begin{tabular}{l}
\leansmall{HasSpectralMassGap.gap\_le\_of\_pos\_spectrum}\\
\leansmall{HasSpectralMassGap.spectrum\_subset\_nonpos\_or\_gap}
\end{tabular}
\end{center}
Thus a positive spectral value of a Hamiltonian with gap \(\Delta\) must lie at
least at \(\Delta\), and the whole spectrum is confined to
\((-\infty,0]\cup[\Delta,\infty)\).  The companion lemma
\[
  \lean{HasSpectralMassGap.not\_spectrum\_of\_pos\_lt}
\]
gives the direct operational exclusion of any \(E\in(0,\Delta)\).
Lean now also proves the exact converse characterization
\[
  \lean{hasSpectralMassGap\_iff\_positive\_and\_spectrum\_subset\_nonpos\_or\_gap}.
\]
So, on this bounded-operator audit surface, a positive \(\Delta\) is a spectral
mass gap exactly when the spectrum lies in \((-\infty,0]\cup[\Delta,\infty)\).

When combined with positive energy, the gap splits the spectrum:
\[
  \lean{HasSpectralMassGap.eq\_zero\_or\_gap\_le\_of\_positive\_energy}.
\]
Every spectral value is either \(0\) or at least \(\Delta\).  This cleanly
separates the vacuum-energy condition from the low-excitation exclusion.
There are also theorem-level forms for non-vacuum spectral values and for the
whole positive-energy spectrum:
\begin{center}
\begin{tabular}{l}
\leansmall{HasSpectralMassGap.gap\_le\_of\_positive\_energy\_of\_ne\_zero}\\
\leansmall{HasSpectralMassGap.spectrum\_subset\_vacuum\_or\_gap\_of\_positive\_energy}
\end{tabular}
\end{center}
The same positive-energy split is now available as an exact iff:
\[
  \lean{hasSpectralMassGap\_iff\_positive\_and\_spectrum\_subset\_vacuum\_or\_gap\_of\_positive\_energy}.
\]
Thus, once positive energy is fixed, the mass-gap predicate is interchangeable
with the statement that every spectral value is either the vacuum value \(0\)
or lies at least at \(\Delta\).

The first-positive-spectral-value interface packages the expected “first
excitation” target without building the underlying Yang--Mills Hamiltonian.
If \(m\) is the first positive spectral value, then lower positive spectral
values are excluded and the whole spectrum lies in
\((-\infty,0]\cup[m,\infty)\):
\begin{center}
\begin{tabular}{l}
\leansmall{HasFirstPositiveSpectralValue.unique}\\
\leansmall{HasFirstPositiveSpectralValue.not\_spectrum\_of\_pos\_lt}\\
\leansmall{HasFirstPositiveSpectralValue.spectrum\_subset\_nonpos\_or\_first}\\
\leansmall{HasFirstPositiveSpectralValue.spectrum\_subset\_vacuum\_or\_first\_of\_positiveEnergy}
\end{tabular}
\end{center}
The uniqueness lemma records that two supplied first positive spectral values
for the same Hamiltonian must be equal.  With positive energy, the nonpositive
side of the split tightens to the vacuum alternative \(E=0\).
The same predicate is also equivalent to the order-theoretic statement that
\(m\) is the least element of the positive spectrum:
\begin{center}
\begin{tabular}{l}
\leansmall{HasFirstPositiveSpectralValue.isLeast\_positiveSpectrumSet}\\
\leansmall{hasFirstPositiveSpectralValue\_iff\_isLeast\_positiveSpectrumSet}\\
\leansmall{HasFirstPositiveSpectralValue.positiveSpectralInf\_eq}
\end{tabular}
\end{center}
This gives the parallel audit value
\(\lean{PositiveSpectralInf}\,A=m\) whenever the first positive spectral value
exists.

The bridge back to the gap predicate is exact:
\[
  \lean{hasFirstPositiveSpectralValue\_iff\_spectrum\_mem\_and\_hasSpectralMassGap}.
\]
Equivalently, a candidate \(m\) is the first positive spectral value precisely
when \(m\) is actually in the spectrum and \(m\) itself is an admitted spectral
mass gap.  The forward direction is also available as
\[
  \lean{HasSpectralMassGap.hasFirstPositiveSpectralValue\_of\_spectrum\_mem}.
\]
The same assumption exactly identifies the admitted spectral gaps:
\begin{center}
\begin{tabular}{l}
\leansmall{HasFirstPositiveSpectralValue.hasSpectralMassGap\_of\_le}\\
\leansmall{HasFirstPositiveSpectralValue.hasSpectralMassGap\_iff}\\
\leansmall{HasFirstPositiveSpectralValue.spectralGapSet\_eq\_Ioc}
\end{tabular}
\end{center}
Thus the bounded audit layer proves
\(\lean{SpectralGapSet}\,A=(0,m]\).  It also supplies finite spectral mass and
collapses the supremum audit to the first excitation:
\begin{center}
\begin{tabular}{l}
\leansmall{HasFirstPositiveSpectralValue.hasFiniteSpectralMass}\\
\leansmall{HasFirstPositiveSpectralValue.spectralGapSup\_eq}\\
\leansmall{HasFirstPositiveSpectralValue.spectralGapSup\_eq\_positiveSpectralInf}
\end{tabular}
\end{center}
These are not Yang--Mills construction theorems; they are downstream checks
that any proposed construction of a Hamiltonian with a first positive spectral
value must satisfy.

The finite-mass audit predicate is linked to actual positive spectrum by
\[
  \lean{hasFiniteSpectralMass\_of\_pos\_spectrum}.
\]
A single positive spectral value gives an upper bound for every admissible gap.
The boundedness surface is now exposed directly as
\[
  \lean{hasFiniteSpectralMass\_iff\_bddAbove\_spectralGapSet}.
\]
Thus the finite-mass predicate is exactly the assertion that the admitted gap
set is bounded above, with the definition's positive bound recovered by taking
the maximum of any upper bound and \(1\).  The new supremum audit value
\[
  \lean{SpectralGapSup}
\]
is deliberately paired with the hypotheses that make it meaningful.  If the
gap set is bounded above and nonempty, then
\[
  \lean{spectralGapSup\_isLUB\_of\_hasFiniteSpectralMass\_of\_gap}
\]
proves that \lean{SpectralGapSup} is the least upper bound of all admitted
gaps.  The pointwise lower and positivity forms are
\begin{center}
\begin{tabular}{l}
\leansmall{HasSpectralMassGap.le\_spectralGapSup\_of\_hasFiniteSpectralMass}\\
\leansmall{spectralGapSup\_pos\_of\_hasFiniteSpectralMass\_of\_gap}
\end{tabular}
\end{center}
and a positive spectral value bounds the supremum by
\[
  \lean{spectralGapSup\_le\_of\_pos\_spectrum\_of\_gap}.
\]
This is the supremum-level audit version of the elementary fact that every
admitted gap must lie below any positive spectral value.  Conversely,
\[
  \lean{not\_hasFiniteSpectralMass\_of\_arbitrarily\_large\_gaps}
\]
shows that admitted gaps above every proposed bound rule out finite spectral
mass.  The earlier special case
\[
  \lean{not\_hasFiniteSpectralMass\_of\_all\_positive\_gaps}
\]
records that if every positive window is accepted as a gap, then no finite
upper bound on gaps can exist.  This distinguishes a genuine finite first
excitation from the degenerate situation where the bounded audit surface sees
no positive spectrum at all.

Finally, the quadratic-form gap has the same downward closure:
\[
  \lean{HasQuadraticMassGap.mono}.
\]
This proof uses only nonnegativity of \(\langle \psi,\psi\rangle\), so it is a
stable preliminary lemma before any later unbounded-operator or spectral theorem
bridge is added.

\section{Exponential Clustering Audit}

The next module is
\[
  \lean{Mettapedia/QuantumTheory/YangMills/Clustering.lean}.
\]
It isolates the Clay clustering consequence without pretending to prove the
analytic mass-gap-to-clustering theorem.  A two-point spatial correlation is
represented abstractly as
\[
  \lean{SpatialCorrelation}=\lean{Space}\to\lean{Space}\to\mathbb{R}.
\]
For such a correlation kernel, the module defines:
\begin{itemize}[leftmargin=1.5em]
  \item \lean{HasExponentialClusteringAtRate}, the fixed-rate estimate
    \(|C(x,y)|\le \exp(-c\,d(x,y))\) beyond some radius;
  \item \lean{HasExponentialClusteringUpTo}, the Clay-style all-rates-below-gap
    statement, requiring the fixed-rate estimate for every \(0<c<\Delta\);
  \item \lean{HasArbitrarilyLargeClusteringViolation}, a concrete obstruction
    saying that every proposed cutoff radius has a farther counterexample to
    the fixed-rate exponential bound;
  \item \lean{SpectralGapClusteringBridge}, an abstract bridge asserting that
    every spectral gap for a Hamiltonian implies clustering of the supplied
    correlation kernel up to that gap.
\end{itemize}

The basic monotonicity and projection lemmas are now formalized:
\begin{center}
\begin{tabular}{l}
\leansmall{HasExponentialClusteringUpTo.at\_rate}\\
\leansmall{HasExponentialClusteringUpTo.mono}\\
\leansmall{SpectralGapClusteringBridge.clusteringUpTo\_of\_gap}\\
\leansmall{SpectralGapClusteringBridge.clusteringAtRate\_of\_gap}
\end{tabular}
\end{center}
Thus, once a proof supplies a spectral-gap-to-clustering bridge, every proposed
gap \(\Delta\) produces each concrete clustering estimate with \(0<c<\Delta\).
This deliberately separates the analytic theorem from the bookkeeping that
future route audits will need.

The first crux-style consequence is also formalized:
\begin{center}
\begin{tabular}{l}
\leansmall{not\_hasExponentialClusteringAtRate\_of\_arbitrarilyLarge\_violation}\\
\leansmall{not\_hasExponentialClusteringUpTo\_of\_rate\_failure}\\
\leansmall{not\_hasSpectralMassGap\_of\_clustering\_rate\_failure}\\
\leansmall{not\_hasSpectralMassGap\_of\_arbitrarilyLarge\_clustering\_violation}\\
\leansmall{spectralGap\_le\_of\_clustering\_rate\_failure}\\
\leansmall{spectralGap\_le\_of\_arbitrarilyLarge\_clustering\_violation}\\
\leansmall{spectralGapSup\_le\_of\_clustering\_rate\_failure\_of\_gap}\\
\leansmall{not\_exists\_hasSpectralMassGap\_of\_all\_positive\_clustering\_rate\_failure}\\
\leansmall{not\_exists\_hasSpectralMassGap\_of\_all\_positive\_arbitrarilyLarge\_clustering\_violation}\\
\leansmall{HasPositiveLongRangeLowerBound.arbitrarilyLargeClusteringViolation}\\
\leansmall{not\_exists\_hasSpectralMassGap\_of\_positiveLongRangeLowerBound}
\end{tabular}
\end{center}
Under a supplied \lean{SpectralGapClusteringBridge}, failure of the exponential
clustering estimate at any rate \(0<c<\Delta\) rules out \(\Delta\) as a
spectral mass gap for that Hamiltonian.  This is not a global Yang--Mills
blocker; it is a reusable audit surface for the manuscript's route from
polymer/tree decay to exponential clustering and then to the OS Hamiltonian
gap.  Any later proof path that claims such a bridge must now provide the
fixed-rate clustering estimates explicitly, or it exposes a concrete rate-level
counterexample target.

The same obstruction is now available in quantitative threshold form.  A
single failed positive clustering rate \(c\) upper-bounds every bridged
spectral gap:
\[
  \lean{HasSpectralMassGap}\;A\;\Delta
  \quad\Longrightarrow\quad \Delta\le c.
\]
If at least one gap is admitted, the audited supremum \(\lean{SpectralGapSup}\)
is also bounded by \(c\).  Finally, if the supplied correlation fails the
fixed-rate exponential estimate at every positive rate, Lean proves that there
is no positive spectral gap at all for that bridged route.  This converts a
family of long-range correlation counterexamples into a direct no-gap
criterion on the abstract Hamiltonian surface.

The newest concrete route-facing form is
\[
  \lean{HasPositiveLongRangeLowerBound}\;\mathrm{corr}\;\varepsilon.
\]
It says that for some \(\varepsilon>0\), beyond every radius there are two
points whose absolute correlation is still at least \(\varepsilon\).  Lean
proves that this lower bound beats every positive exponential decay rate:
\[
  \lean{HasPositiveLongRangeLowerBound.arbitrarilyLargeClusteringViolation}.
\]
Combining this with the spectral-gap-to-clustering bridge gives
\[
  \lean{not\_exists\_hasSpectralMassGap\_of\_positiveLongRangeLowerBound}.
\]
Thus a persistent nonzero long-range tail in the proposed two-point function is
not merely a missing estimate; on this bridge surface it directly rules out
all positive spectral gaps.  This provides a concrete target for Yang--Mills
route audits: either prove genuine exponential decay of the centered
correlation functions, or avoid any lower-bound tail that persists uniformly at
arbitrarily large separations.

\section{RG Bootstrap Audit}

The next route-facing module is
\[
  \lean{Mettapedia/QuantumTheory/YangMills/RGBootstrap.lean}.
\]
This file targets the manuscript's multiscale RG contraction claim in a narrow,
checkable way.  It does not prove the polymer RG step.  Instead, it isolates the
finite arithmetic and recurrence facts that any later formalized RG proof must
instantiate.

The module defines
\[
  \lean{irrelevantScale}\;b\;d_{\max}=b^{3-d_{\max}}
\]
using an integer exponent, together with the predicates
\lean{HasRGContraction} and \lean{HasExtendedExtractionContraction}.  It also
defines \lean{AffineRGStepBound}, the recurrence surface
\[
  a_{n+1}\le \kappa a_n+\varepsilon.
\]
The basic finite bootstrap lemma is
\[
  \lean{AffineRGStepBound.bound\_by\_invariant\_interval}.
\]
If \(a_0\le B\), \(\kappa\ge 0\), and
\(\varepsilon\le(1-\kappa)B\), then every finite iterate satisfies
\(a_n\le B\).  This is the algebraic core of a smallness bootstrap once the
analytic RG estimates have supplied a one-step recurrence.

The same module now also exposes the exact finite-iterate envelope
\[
  \lean{affineRGEnvelope}\;\kappa\;\varepsilon\;a_0,
\]
with theorem
\[
  \lean{AffineRGStepBound.le\_affineRGEnvelope}.
\]
Thus any sequence satisfying the one-step affine bound is dominated at every
finite scale by the recurrence \(x_{n+1}=\kappa x_n+\varepsilon\) started at
\(a_0\).  In the homogeneous case, Lean proves
\[
  \lean{affineRGEnvelope\_zero\_error}
\]
and the resulting wrapper
\[
  \lean{AffineRGStepBound.bound\_by\_geometric\_of\_zero\_error},
\]
so the envelope reduces exactly to \(x_n=\kappa^n a_0\).  This keeps the
finite-scale bookkeeping proved independently of the still-open analytic RG
step estimate.

The numerical extended-extraction audit is exact:
\begin{center}
\begin{tabular}{l}
\leansmall{irrelevantScale\_two\_sixteen}\\
\leansmall{irrelevantScale\_two\_four}\\
\leansmall{irrelevantScale\_two\_fourteen}\\
\leansmall{irrelevantScale\_two\_fifteen}\\
\leansmall{irrelevantScale\_eq\_inv\_pow\_of\_three\_le}\\
\leansmall{rgContraction\_iff\_constant\_lt\_pow\_of\_pos\_of\_three\_le}\\
\leansmall{rgContraction\_forces\_constant\_lt\_pow\_of\_pos\_of\_three\_le}\\
\leansmall{not\_rgContraction\_2224\_of\_pow\_le\_of\_pos\_of\_three\_le}\\
\leansmall{irrelevantScale\_two\_eq\_inv\_pow\_of\_three\_le}\\
\leansmall{rgContraction\_two\_iff\_constant\_lt\_pow\_of\_three\_le}\\
\leansmall{rgContraction\_two\_forces\_constant\_lt\_pow\_of\_three\_le}\\
\leansmall{rgContraction\_2224\_two\_sixteen}\\
\leansmall{rgContraction\_2224\_two\_fifteen}\\
\leansmall{rgGain\_2224\_two\_sixteen\_eq}\\
\leansmall{rgGain\_2224\_two\_sixteen\_lt\_one\_third}\\
\leansmall{not\_rgContraction\_2224\_two\_four}\\
\leansmall{rgGain\_2224\_two\_four\_eq}\\
\leansmall{rgGain\_2224\_two\_fourteen\_eq}\\
\leansmall{rgGain\_2224\_two\_fifteen\_eq}\\
\leansmall{rgContraction\_two\_fourteen\_forces\_constant\_lt\_2048}\\
\leansmall{not\_rgContraction\_two\_fourteen\_of\_ge\_2048}\\
\leansmall{not\_rgContraction\_two\_of\_ge\_2048\_of\_dmax\_le\_fourteen}\\
\leansmall{not\_rgContraction\_2224\_two\_fourteen}\\
\leansmall{not\_rgContraction\_2224\_two\_of\_dmax\_le\_fourteen}\\
\leansmall{rgContraction\_of\_nonneg\_le\_2224\_two\_sixteen}
\end{tabular}
\end{center}
Lean proves
\[
  \lean{irrelevantScale}\;2\;16=2^{-13}=1/8192
\]
and the exact gain
\[
  2224\cdot 2^{-13}=139/512<1/3.
\]
It now also proves the general scale formula for every \(d_{\max}\ge 3\):
\[
  \lean{irrelevantScale}\;b\;d_{\max}=b^{-(d_{\max}-3)}
    =\bigl(b^{d_{\max}-3}\bigr)^{-1}.
\]
For every positive natural block factor \(b\), Lean proves the exact threshold
iff
\[
  \lean{HasExtendedExtractionContraction}\;C\;b\;d_{\max}
  \quad\Longleftrightarrow\quad
  0\le C\ \wedge\ C<b^{d_{\max}-3}.
\]
Consequently, changing the block factor does not evade the audit; it merely
changes the required explicit power threshold.  In particular,
\[
  b^{d_{\max}-3}\le 2224
  \quad\Longrightarrow\quad
  \neg\lean{HasExtendedExtractionContraction}\;2224\;b\;d_{\max}
\]
for \(b>0\) and \(d_{\max}\ge3\).  The block-factor-\(2\) specialization is:
\[
  \lean{irrelevantScale}\;2\;d_{\max}=2^{-(d_{\max}-3)}
    =\bigl(2^{d_{\max}-3}\bigr)^{-1},
\]
and the exact iff
\[
  \lean{HasExtendedExtractionContraction}\;C\;2\;d_{\max}
  \quad\Longleftrightarrow\quad
  0\le C\ \wedge\ C<2^{d_{\max}-3}.
\]
Thus any block-factor-\(2\) contraction certificate at depth
\(d_{\max}\ge 3\) must place its recombination constant below the explicit
power threshold \(2^{d_{\max}-3}\).  The earlier \(d_{\max}=14\) statement
\(C<2048\) is now just the nearest-even-depth specialization of this exact
threshold, not an isolated arithmetic calculation.
It also proves that the same constant fails to contract at the naive
\(d_{\max}=4\) scale, where the factor is \(1/2\) and the exact gain is
\(1112\).  The nearest even lower extraction depth also fails at the advertised
constant:
\[
  \lean{irrelevantScale}\;2\;14=2^{-11}=1/2048,\qquad
  2224\cdot 2^{-11}=139/128>1.
\]
Moreover, any contraction certificate at \(d_{\max}=14\) forces the linear
constant to be strictly below \(2048\), so every proposed \(d_{\max}=14\)
certificate with \(C\ge 2048\) already fails.  Lean then packages the same
threshold uniformly for every lower extraction depth \(d_{\max}\le 14\) with
block factor \(2\):
\[
  2048 \le C
  \quad\Longrightarrow\quad
  \neg \lean{HasExtendedExtractionContraction}\;C\;2\;d_{\max}.
\]
So the lower-depth repair burden is no longer merely ``beat \(2224\)''; any
block-factor-\(2\) repair below depth \(15\) must actually sharpen the
recombination constant below the exact threshold \(2048\).  To avoid
overclaiming, the audit also records that the raw natural-number cutoff
\(d_{\max}=15\) has
\[
  \lean{irrelevantScale}\;2\;15=2^{-12}=1/4096,\qquad
  2224\cdot 2^{-12}=139/256<1,
\]
so the same advertised constant contracts there.  The route-level crux is
therefore an even canonical-dimension obstruction, not a claim that every
natural cutoff below \(16\) fails.  This is not a global Yang--Mills blocker,
but it is a precise
same-constant obstruction: any repair that keeps the manuscript's advertised
combinatorial bound while dropping the even extended extraction depth below
\(16\) must replace this contraction argument with something genuinely
different or justify a nonstandard odd cutoff.
That residual odd-cutoff loophole is now audited separately in
\[
  \lean{Mettapedia/QuantumTheory/YangMills/ExtractionOddCutoff.lean}.
\]
It defines \lean{VanishesOnOddDimensions} and proves
\[
  \lean{extendedExtraction\_fifteen\_eq\_fourteen\_of\_vanishOnOddDimensions}
\]
and
\[
  \lean{higherExtractionBlock\_fifteen\_eq\_fourteen\_of\_vanishOnOddDimensions}.
\]
So on the manuscript's current even-dimensional extracted-operator surface, the
putative odd repair \(d_{\max}=15\) is algebraically identical to
\(d_{\max}=14\): it adds no new extracted couplings and no new higher block.
Therefore the raw natural-number contraction at \(15\) is not by itself a live
repair on that surface; using it would first require genuinely new odd
canonical-dimension operators absent from the present extraction model.
The RG-facing remainder object is now formalized separately in
\[
  \lean{Mettapedia/QuantumTheory/YangMills/ExtractionRemainder.lean}.
\]
This file defines
\[
  \lean{extractionRemainder}\;d\;\mathrm{coeff}
    = \mathrm{coeff} - \lean{extendedExtraction}\;d\;\mathrm{coeff},
\]
proves the pointwise split of coefficients into extracted part plus remainder,
and shows that the remainder vanishes exactly when the coefficients are
supported up to the cutoff.  Most importantly for the odd-cutoff loophole, it
proves
\[
  \lean{extractionRemainder\_fifteen\_eq\_fourteen\_of\_vanishOnOddDimensions}.
\]
So on the manuscript's even-dimensional surface, not only the extracted
couplings but the actual post-extraction remainder fed into the RG bootstrap is
unchanged when one replaces cutoff \(14\) by \(15\).  The raw \(2^{-12}\)
arithmetic therefore does not correspond to a genuinely different remainder
being bootstrapped on this surface.
That collapse is now pushed to the actual route-witness level in
\[
  \lean{Mettapedia/QuantumTheory/YangMills/ExtractionStateRouteCollapse.lean}.
\]
This file defines the downstream RG state
\[
  \lean{ExtractionState}\;\alpha
    := (\lean{extendedExtraction}\;d\;\mathrm{coeff},
        \lean{extractionRemainder}\;d\;\mathrm{coeff}),
\]
proves
\[
  \lean{extractionState\_fifteen\_eq\_fourteen\_of\_vanishOnOddDimensions},
\]
and then packages the route-facing abstraction
\[
  \lean{DeterminedByExtractionState}.
\]
Its manuscript-facing corollaries are
\[
  \lean{extractionStateRouteWitness\_fifteen\_iff\_fourteen\_of\_vanishOnOddDimensions}
\]
and
\[
  \lean{not\_extractionStateRouteWitness\_fifteen\_of\_not\_fourteen\_on\_oddVanishingSurface}.
\]
So once a proposed odd-cutoff repair witness depends only on the extracted
local terms together with the post-extraction remainder, Lean reduces the
\(d_{\max}=15\) presentation exactly to the already-blocked \(d_{\max}=14\)
surface.  The remaining burden is no longer “show that \(15\) contracts
arithmetically”; it is to exhibit some genuinely new downstream datum that is
not already contained in this extraction-state pair.
That bridge is now pushed all the way back to the existing RG route crux in
\[
  \lean{Mettapedia/QuantumTheory/YangMills/ExtractionStateRGCrux.lean}.
\]
This file defines
\[
  \lean{ForcesSameConstantFourteenContraction}
\]
for extraction-state predicates \(P\): every cutoff-\(14\) witness of \(P\)
already yields the manuscript's advertised same-constant contraction
\(\lean{HasExtendedExtractionContraction}\;2224\;2\;14\).
With that compatibility premise in hand, Lean proves
\[
  \lean{sameConstantFourteenGapRoute\_of\_oddVanishingExtractionStateGapRoute},
\]
which collapses any odd-vanishing cutoff-\(15\) extraction-state route to the
pointed blocked \(d_{\max}=14\) gap surface, and then proves
\[
  \lean{sameConstantLowerExtractionGapRouteCertificate\_of\_oddVanishingExtractionStateGapRoute}
\]
followed by the crux-out theorem
\[
  \lean{not\_oddVanishingExtractionStateGapRoute\_of\_forcesSameConstantFourteenContraction}.
\]
So the exact missing compatibility lemma is now isolated in Lean: if an
alleged odd-cutoff repair extracts only the same local terms and remainder, and
those data already force the manuscript's same-constant \(d_{\max}=14\)
contraction, then the whole \(d_{\max}=15\) route is formally impossible.

The manuscript also makes a separate algebraic coherence claim about the two
extraction schemes:
\[
  P_{\le 4}=P_{\le 4}\circ P_{\le d_{\max}}.
\]
This is now formalized in
\[
  \lean{Mettapedia/QuantumTheory/YangMills/ExtractionProjection.lean}.
\]
The file models a local expansion only at the canonical-dimension coefficient
level, via \lean{extractLE}.  It does not yet define the actual jet-based local
functional projector from the manuscript.  On this honest algebraic surface,
Lean proves:
\begin{center}
\begin{tabular}{l}
\leansmall{extractLE\_idem}\\
\leansmall{extractLE\_compose}\\
\leansmall{extractLE\_compose\_of\_le}\\
\leansmall{extractLE\_eq\_self\_iff\_supportedUpTo}\\
\leansmall{standardExtraction\_compose\_extendedExtraction\_of\_four\_le}
\end{tabular}
\end{center}
Thus the manuscript's compatibility identity is now formally available once the
actual Yang--Mills extraction operator is shown to act by truncation on the
relevant canonical-dimension coefficients.
The same file now also isolates the extra extracted couplings above the standard
cutoff as \lean{higherExtractionBlock}.  On the same coefficient-truncation
surface, Lean proves
\[
  \lean{standardExtraction\_higherExtractionBlock\_eq\_zero}
\]
and
\[
  \lean{extendedExtraction\_eq\_standard\_plus\_higherExtractionBlock\_of\_four\_le}.
\]
So for \(d_{\max}\ge 4\), the extended extraction splits exactly into the
standard \(P_{\le 4}\) part plus the higher \(5,\ldots,d_{\max}\) block, and
that higher block is invisible to the standard extraction.  This is still not a
construction of the manuscript's jet-based projector, but it sharpens the exact
algebra any such projector must satisfy.
The same surface is now also explicitly additive:
\[
  \lean{extractLE\_zero},
  \qquad
  \lean{extractLE\_add}.
\]
Consequently Lean proves
\[
  \lean{standardExtraction\_add\_higherExtractionBlock\_eq}.
\]
So adding a pure higher extracted block to any coefficient family cannot change
its \(P_{\le 4}\) projection.  At the manuscript level, this means any claimed
effect of the extra \(5,\ldots,d_{\max}\) couplings on the standard
extraction/beta-drift channel must come from additional structure beyond this
bare truncation algebra.  The same coefficient surface now also identifies that
higher block exactly as the residual
\[
  \lean{higherExtractionBlock\_eq\_extendedExtraction\_sub\_standard\_of\_four\_le},
\]
and Lean proves the equivalent vanishing statement
\[
  \lean{standardExtraction\_extendedExtraction\_sub\_standard\_eq\_zero\_of\_four\_le}.
\]
So after subtracting the standard extraction from the extended one, the
remaining \(5,\ldots,d_{\max}\) residue still has zero \(P_{\le 4}\)
projection.
The downstream consequence is now packaged separately in
\[
  \lean{Mettapedia/QuantumTheory/YangMills/ExtractionObservable.lean}.
\]
It defines observables \(\Phi\) that are determined by the standard
extraction, in the sense that equal \(P_{\le 4}\) projections force equal
\(\Phi\)-values.  On that surface Lean proves
\[
  \lean{DeterminedByStandardExtraction.extendedExtraction\_eq},
  \qquad
  \lean{DeterminedByStandardExtraction.add\_higherExtractionBlock\_eq},
\]
and
\[
  \lean{DeterminedByStandardExtraction.extendedExtraction\_sub\_standard\_eq\_zero}.
\]
So any beta-drift observable or renormalized coupling functional that genuinely
depends only on \(P_{\le 4}\) is exactly unchanged by replacing a coefficient
family with its extended extraction, by adding the pure higher
\(5,\ldots,d_{\max}\) block, or by passing to the residual
\(P_{\le d_{\max}} - P_{\le 4}\).  This makes the route burden explicit: if the
manuscript wants the extra extracted couplings to affect the beta channel, that
effect has to enter through structure beyond bare dependence on \(P_{\le 4}\).
The same file now also isolates the next route-relevant step: if a candidate
beta observable \(G\) decomposes as a standard-extraction contribution plus an
error term \(E\), then all higher-coupling sensitivity is forced into \(E\).
Lean proves
\[
  \lean{difference\_eq\_errorDifference\_of\_add\_higherExtractionBlock},
  \qquad
  \lean{difference\_eq\_errorDifference\_of\_extendedExtraction},
\]
and
\[
  \lean{difference\_from\_zero\_eq\_errorDifference\_of\_extendedExtraction\_sub\_standard}.
\]
So under the manuscript's own compatibility premise, any claimed effect of the
extra extracted couplings on beta drift must be carried entirely by whatever
additional higher-order mixing/error term is supplied; it cannot come from the
bare \(P_{\le 4}\) contribution itself.
This now reaches the manuscript's Remark G.2 language about ``the same error
bounds.''  The same file proves
\[
  \lean{norm\_difference\_le\_of\_add\_higherExtractionBlock\_errorBound},
  \qquad
  \lean{norm\_difference\_le\_of\_extendedExtraction\_errorBound},
\]
and
\[
  \lean{norm\_difference\_le\_of\_extendedExtraction\_sub\_standard\_errorBound}.
\]
So if the higher-coupling contribution is bounded inside an error term \(E\)
with some norm estimate, then the full beta-like observable \(G\) inherits
exactly the same bound.  In particular, any advertised \(O(\beta_j^{-2})\)
control has to be proved at the error-term layer; the \(P_{\le 4}\) part adds
no extra dependence on the higher extracted couplings.
The same file now also packages the genuinely stepwise drift statement:
\[
  \lean{stepDifference\_eq\_errorStepDifference\_of\_extendedExtraction}
\]
and its norm corollary
\[
  \lean{norm\_stepDifference\_le\_of\_extendedExtraction\_errorBounds}.
\]
So when one compares the beta-drift increment between two RG steps before and
after passing to extended extraction, the entire perturbation is exactly the
corresponding perturbation of the error term, and any bound on those two
error-term responses transfers directly to the drift increment.
The same surface now turns that observation into explicit threshold blockers:
\begin{center}
\begin{tabular}{l}
\leansmall{abs\_difference\_le\_of\_extendedExtraction\_errorBound}\\
\leansmall{extendedExtraction\_le\_target\_of\_le\_target\_sub\_of\_errorBound}\\
\leansmall{extendedExtraction\_lt\_target\_of\_lt\_target\_sub\_of\_errorBound}\\
\leansmall{stepDifference\_le\_target\_of\_le\_target\_sub\_of\_errorBounds}\\
\leansmall{stepDifference\_lt\_target\_of\_lt\_target\_sub\_of\_errorBounds}
\end{tabular}
\end{center}
These theorems formalize the route consequence of the manuscript's ``same error
bounds'' language.  If a \(P_{\le 4}\)-determined real-valued observable is
already at least \(B\) below a target \(T\), and the extended-extraction error
response is bounded by \(B\), then the extended-extracted observable still lies
at or below \(T\).  Likewise, if a stepwise beta-drift inequality already
fails by a margin at least \(B_{\mathrm{next}} + B_{\mathrm{cur}}\), then
extended extraction cannot repair that failure unless the combined error-term
response exceeds that same margin.  So a repair route cannot simply reuse the
old \(O(\beta_j^{-2})\) error budget and claim a threshold or sign correction
for free; it must exhibit genuinely larger error response exactly where the
standard \(P_{\le 4}\) observable misses the required target.

\section{Same-Constant Lower-Extraction Crux}

The first dedicated crux-out module is
\[
  \lean{Mettapedia/QuantumTheory/YangMills/RGCrux.lean}.
\]
It packages the lower-extraction repair attempt as
\[
  \lean{SameConstantLowerExtractionRGCertificate}\;d_{\max}.
\]
This certificate means \(d_{\max}\le 14\) together with the same manuscript
linear contraction certificate
\[
  \lean{HasExtendedExtractionContraction}\;2224\;2\;d_{\max}.
\]
Lean proves
\[
  \lean{not\_sameConstantLowerExtractionRGCertificate}
\]
and the quantified version
\[
  \lean{not\_exists\_sameConstantLowerExtractionRGCertificate}.
\]
Thus there is no same-bound RG certificate at any extraction depth
\(d_{\max}\le 14\).

The module also packages the manuscript-shaped version
\[
  \lean{SameConstantEvenBelowSixteenRGCertificate}\;d_{\max},
\]
which requires \(d_{\max}<16\), \(d_{\max}\) even, and the same contraction
certificate.  The arithmetic lemma
\[
  \lean{even\_lt\_sixteen\_le\_fourteen}
\]
proves that such a cutoff is at most \(14\), and hence Lean proves
\[
  \lean{not\_sameConstantEvenBelowSixteenRGCertificate}
\]
and
\[
  \lean{not\_exists\_sameConstantEvenBelowSixteenRGCertificate}.
\]
This is the preferred statement for comparison with the manuscript's canonical
dimensions \(4,6,\ldots,d_{\max}\).

The same file now packages the threshold-level variant
\[
  \lean{AtLeast2048EvenBelowSixteenRGCertificate}\;C\;d_{\max},
\]
which requires \(d_{\max}<16\), \(d_{\max}\) even, \(2048\le C\), and the
same block-factor-\(2\) contraction certificate.  Lean proves
\[
  \lean{not\_atLeast2048EvenBelowSixteenRGCertificate}
\]
and
\[
  \lean{not\_exists\_atLeast2048EvenBelowSixteenRGCertificate}.
\]
So the obstruction is no longer tied only to the manuscript's literal
\(2224\): any lower-even extraction repair that keeps the recombination
constant at or above \(2048\) already fails before any clustering or
Osterwalder--Schrader argument is reached.

The same obstruction is lifted to the surrounding mass-gap route certificate:
\[
  \lean{SameConstantLowerExtractionGapRouteCertificate}.
\]
This predicate includes a candidate spectral gap, a spectral-gap-to-clustering
bridge, and a lower-extraction same-constant RG certificate.  Lean proves
\[
  \lean{not\_sameConstantLowerExtractionGapRouteCertificate}
\]
and, for the even manuscript-shaped route certificate,
\[
  \lean{not\_sameConstantEvenBelowSixteenGapRouteCertificate}.
\]
The strengthened threshold version is also lifted as
\[
  \lean{AtLeast2048EvenBelowSixteenGapRouteCertificate}
\]
with theorem
\[
  \lean{not\_atLeast2048EvenBelowSixteenGapRouteCertificate}.
\]
and the pointed \(d_{\max}=14\) form
\[
  \lean{not\_sameConstantFourteenExtractionGapRouteCertificate}.
\]
This theorem deliberately does not refute a Yang--Mills mass gap.  It refutes
one repair corridor for the manuscript route: keeping the advertised constant
\(2224\), using block factor \(2\), and lowering the extended extraction depth
below \(16\).  A viable repair must either sharpen the \(d_{\max}=14\)
constant below \(2048\), keep \(d_{\max}=16\), or replace the linear
contraction certificate with a genuinely different argument.

\section{Regression}

The companion module
\[
  \lean{Mettapedia/QuantumTheory/YangMills/MassGapRegression.lean}
\]
contains theorem-level wrappers for the positive-gap projection, pointwise
spectral obstruction, downward closure, positive-energy spectral splitting, and
quadratic-form monotonicity.  It also checks the threshold/subset forms of the
spectral exclusion, including the positive-energy non-vacuum threshold and the
finite-mass consequences, including the bounded-above characterization, the
least-upper-bound API for \lean{SpectralGapSup}, its positivity consequence,
the positive-spectrum upper bound, the first-positive-spectral-value exact
endpoint equivalence, uniqueness, gap-set/supremum identifications, and the
positive-spectrum infimum comparison, and the unbounded-gap obstruction.

The companion clustering regression module
\[
  \lean{Mettapedia/QuantumTheory/YangMills/ClusteringRegression.lean}
\]
pins the fixed-rate projection, monotonicity in the endpoint \(\Delta\), the
rate-level violation blocker, the bridge projections from spectral gap to
clustering, the final contrapositive obstruction from a clustering-rate
failure back to failure of the proposed spectral mass gap, the quantitative
gap upper-bound form, the supremum upper-bound form, and the all-positive-rates
no-gap criterion.  It also pins the positive-long-range-lower-bound route
blocker and its conversion into persistent fixed-rate clustering violations.

The RG bootstrap regression module
\[
  \lean{Mettapedia/QuantumTheory/YangMills/RGBootstrapRegression.lean}
\]
pins the exact irrelevant-scale values for \(d_{\max}=4,14,15,16\), the
general reciprocal-power threshold for every \(b\) and \(d_{\max}\ge3\), the
exact positive-block-factor contraction iff \(0\le C\) and
\(C<b^{d_{\max}-3}\), the advertised-constant blocker when
\(b^{d_{\max}-3}\le2224\), the block-factor-\(2\) specialization, the
advertised extended-extraction contraction, the raw
\(d_{\max}=15\) contraction, the exact gain values \(139/512<1/3\),
\(139/256\), \(139/128\), and \(1112\), the naive \(d_{\max}=4\)
noncontraction, the
nearest-lower-even-depth \(d_{\max}=14\) noncontraction and threshold
\(C<2048\), the stronger blocker for every \(C\ge 2048\) at
\(d_{\max}=14\) and, more generally, at every \(d_{\max}\le 14\), the monotone
variant for any nonnegative constant bounded by \(2224\), the
invariant-interval bootstrap lemma, and the finite-envelope/geometric-zero-error
wrappers.

The extraction observable regression module
\[
  \lean{Mettapedia/QuantumTheory/YangMills/ExtractionObservableRegression.lean}
\]
pins the factorization criterion through \(P_{\le 4}\), the invariance under
added higher extracted blocks and extended extraction, the vanishing of the
extended-minus-standard residue, the pure-error identities for higher-block,
extended-extraction, and residual responses, the inherited norm bounds from the
error term, and the step-difference/error-difference identity together with its
norm bound.  It now also pins the real-valued threshold forms: absolute
observable perturbation is bounded by the same error budget, any observable
already below a target by more than that budget stays below the target after
extended extraction, and the same quantitative obstruction holds for
step-drift inequalities using the combined two-step error response.

The odd-cutoff regression module
\[
  \lean{Mettapedia/QuantumTheory/YangMills/ExtractionOddCutoffRegression.lean}
\]
pins the successor-cutoff collapse when the new coefficient vanishes, the
odd-successor collapse on odd-vanishing coefficient families, and the
manuscript-facing specializations
\(\lean{extendedExtraction}\;15=\lean{extendedExtraction}\;14\) and
\(\lean{higherExtractionBlock}\;15=\lean{higherExtractionBlock}\;14\) on that
even-dimensional surface.

The extraction remainder regression module
\[
  \lean{Mettapedia/QuantumTheory/YangMills/ExtractionRemainderRegression.lean}
\]
pins the coefficient decomposition into extracted part plus remainder, the
pointwise vanishing/agreement behavior below and above the cutoff, the exact
characterization of zero remainder by supportedness up to the cutoff, and the
manuscript-facing equality of the actual remainders at cutoffs \(15\) and
\(14\) on odd-vanishing coefficient families.

The extraction-state route-collapse regression module
\[
  \lean{Mettapedia/QuantumTheory/YangMills/ExtractionStateRouteCollapseRegression.lean}
\]
then pins the full extraction-state equality at cutoffs \(15\) and \(14\), the
generic route-witness collapse for predicates factored through that state, a
positive odd-vanishing example where the witness does collapse, and a negative
odd-sensitive example showing that the odd-vanishing hypothesis is genuinely
used.

The extraction-state RG-crux regression module
\[
  \lean{Mettapedia/QuantumTheory/YangMills/ExtractionStateRGCruxRegression.lean}
\]
then pins the route-level collapse into the pointed \(d_{\max}=14\) gap
surface, the induced lower-extraction gap certificate, the generic no-go
theorem, and a concrete negative example showing that not every arbitrary
extraction-state predicate can force the blocked cutoff-\(14\) contraction.

The RG crux regression module
\[
  \lean{Mettapedia/QuantumTheory/YangMills/RGCruxRegression.lean}
\]
pins the nonexistence of same-constant lower-extraction RG certificates, the
nonexistence of such certificates at any \(d_{\max}\le 14\), the even-cutoff
lemma reducing \(d_{\max}<16\) and even to \(d_{\max}\le14\), and the lifted
route-certificate blockers for the surrounding mass-gap/clustering bridge
surface.  It now also pins the threshold-level nonexistence of every
even-cutoff lower-extraction certificate with recombination constant at least
\(2048\), together with the corresponding lifted route blocker.

\section{What Remains}

This layer does not solve the Clay problem.  The hard work still includes:
\begin{itemize}[leftmargin=1.5em]
  \item constructing a nontrivial four-dimensional quantum Yang--Mills theory;
  \item choosing a mathematically adequate gauge group and field-strength
    formalization;
  \item connecting bounded-operator placeholders to the physically relevant
    unbounded Hamiltonian theory;
  \item proving a genuine spectral gap for that Hamiltonian;
  \item relating the spectral gap to clustering and local operator estimates.
  \item instantiating \lean{SpatialCorrelation} with centered Yang--Mills
    local-operator vacuum correlations and proving the polymer/tree-decay
    estimates needed by the manuscript's OS reconstruction route.
  \item instantiating \lean{AffineRGStepBound} with the actual polymer norm
    recurrence, including the one-step RG estimates and all constants feeding
    the advertised \(2224\cdot2^{-13}<1\) contraction.
  \item for any lower-even-extraction repair with \(d_{\max}<16\), either
    sharpening the \(d_{\max}=14\) recombination constant below \(2048\),
    justifying a nonstandard odd cutoff such as \(d_{\max}=15\), or abandoning
    the same linear contraction certificate.
\end{itemize}

The value of the current step is that future claims now have a precise Lean
interface for the mass-gap and clustering audit obligations.

\end{document}
